Paper 1 has the higher mean ($35.75$ rather than $25.75$) and lower standard deviation ($7.79$ rather than $11.91$), so its marks are higher on average and less spread out.

A zero on paper 2 despite a near-average paper-1 mark merits checking. Absence, a recording error, or a genuine poor performance are possible. The summary alone cannot establish that the point is an outlier; inspect the actual scatter and the student's record.

The original totals are $\sum x=572$ and $\sum y=412$. After omission they become $534,412$, so the new means are
\[
\boxed{\bar x=35.6,\qquad\bar y=\dfrac{412}{15}=27.4666\ldots}.
\]
The removed product is $38\cdot0=0$, so $\sum xy=15837$ remains unchanged, and
\[
S_{xy}=15837-\dfrac{534\cdot412}{15}=1169.8.
\]
Thus
\[
\boxed{r=\dfrac{1169.8}{\sqrt{965.6\cdot1561.7}}\approx0.952608}.
\]
The slope is $b=\dfrac{S_{xy}}{S_{xx}}=1.2114747\ldots$ and the intercept is $a=\bar y-b\bar x=-15.6618337\ldots$. Hence the fitted line is
\[
\boxed{y\approx-15.6618+1.21147x}.
\]
The stronger correlation describes the remaining sample, but is not by itself a valid reason to delete a genuine observation. Exclusion requires a defensible reason such as a verified error or a population definition that excludes a missed examination. Substituting $x=38$ into the provisional line gives $\boxed{y\approx30.4}$ marks.
